\section{Discussion}
\label{sec:discussion}

The results support the central hypothesis: sparse target-well measurements are
more useful when they are enforced through a measurement-consistency term over a
learned latent manifold than when they are used only as additional supervised
features. In the cross-well Vc benchmark, the training set is small and the
held-out well has a shifted target distribution. Direct baselines must learn a
complete \([u,b]\to y\) map from few windows. CLP-CSGM Ridge decomposes the task
into two simpler components: learn a plausible target manifold from training
windows, then select a nearby latent code that is consistent with sparse
measurements in the target well.

The ridge prior is important in this decomposition. It provides a stable
log-conditioned center \(z_0(u)\) in regimes where a more flexible predictor
could overfit. The sparse measurements then act as a local correction rather
than as the only source of information. This explains why the method performs
especially well in low-data and low-\(\rho\) settings.

The real-well F03 experiment gives a more nuanced result. CLP-CSGM Ridge is the
best method on average and is strongest when \(\rho\) is small. However, AE
\([u,b]\) slightly outperforms it at larger \(\rho\). This loss is not
inconsistent with the method objective. CLP-CSGM is designed for sparse
assimilation under a generative prior; when many target samples are available,
a direct model can exploit them without paying the cost or bias of staying near
the learned latent manifold.

Several limitations remain. First, the autoencoder is intentionally small, so
representation error may limit reconstruction of unusual geological patterns.
Second, the theoretical guarantees of CSGM rely on assumptions such as suitable
measurement operators and generator ranges; our coordinate subsampling and
field data do not directly satisfy a proven guarantee. Third, latent
optimization is nonconvex and more expensive than direct regression. Finally,
the experiments use a limited number of wells, so broader geological validation
is needed before claiming field-wide generality.

These limitations also clarify the contribution. CLP-CSGM Ridge is not a
replacement for all direct petrophysical regressors. It is a targeted method
for regimes where dense logs provide a useful but incomplete prior and sparse
core measurements should be honored as physical evidence in the target well.
